\begin{papercard}{fereday2006}
\cf{Stages} \stagetag{C} (also \stagetag{I}). A worked example of \emph{hybrid thematic analysis}: a template of a priori, theory-derived codes (Crabtree \& Miller) combined with data-driven inductive coding (Boyatzis), carried out in six documented stages. It is the usual citation for justifying a codebook that starts deductive and then grows.

\cf{Context and method} A doctoral study of performance feedback in nursing, framed by Schutz's social phenomenology. Data: 10 focus groups (clinicians and managers) and 16 policy documents, all coded in NVivo \pc{fereday2006}{4}{Following data collection from 10 focus groups and a document analysis of 16 policies or procedures}. The analysis combines the inductive approach of Boyatzis with the deductive a priori template of Crabtree and Miller, so that theory is integral while themes can still emerge \pc{fereday2006}{4}{This approach complemented the research questions by allowing the tenets of social phenomenology}. The template (codebook) is defined \emph{before} in-depth analysis, here from the research question and framework rather than a scan of the text \pc{fereday2006}{4}{for this study, the template was developed a priori, based on the research question and the theoretical framework}. Data collection and analysis were concurrent and iterative, with earlier stages reread to keep themes grounded \pc{fereday2006}{4}{we reread the previous stages of the process before undertaking further analysis}.

\cf{Rigor framework} Schutz's postulates: logical consistency, subjective interpretation, adequacy. Rigor is shown through transparency of the analytic steps, quotations from raw data, member checks (limited here to end-of-session summaries) and practitioner uptake \pc{fereday2006}{3}{Interpretive rigor requires the researcher to demonstrate clearly how interpretations of the data have been achieved}.

\cf{Six stages (Figure 1)} \textbf{1. Developing the code manual}: a template with six broad categories; each code has a \emph{label, a definition and a description of how to know when the theme occurs} (Boyatzis format, Table 1) \pc{fereday2006}{5}{For this study, codes were written with reference to Boyatzis (1998)}. \textbf{2. Testing reliability of the code}: two documents were coded by the researcher and the supervisor, compared, and no changes to the template were needed \pc{fereday2006}{6}{Two performance appraisal documents from health care organizations were selected as test pieces}. \textbf{3. Summarising data and identifying initial themes}: paraphrase each transcript by trigger question; frequency is not the aim, so a single comment counts as much as a repeated one \pc{fereday2006}{7}{It is important to note that a content analysis was not the aim of the data analysis}. \textbf{4. Applying the template and additional coding}: codes entered as NVivo nodes and applied across all three data sets; guided but not confined by the template; inductive codes added as new or expanded codes. Example: ``trust and respect'' started inside the We-relation code and became a separate data-driven code \pc{fereday2006}{9}{these additional codes were either separate from the predetermined codes or they expanded a code from the manual}. \textbf{5. Connecting the codes and identifying themes}: cluster by research question, compare data sources and groups (Table 6). \textbf{6. Corroborating and legitimating coded themes}: re-check clusters against the earlier stages to avoid unintentionally ``seeing'' expected data, and build first- to second-order themes (Table 7) \pc{fereday2006}{11}{Fabricating evidence can be a common problem in the process of interpreting data}.

\cf{Limitation stated} Coding by one person with the analysis discussed with a supervisor gave consistency but not multiple perspectives. Suggested remedy: several coders, expert panel or participants \pc{fereday2006}{12}{the data were coded and themes identified in the data by one person and the analysis then discussed with a supervisor}.

\cf{Use for this thesis} A citable model for the \emph{codebook origin story}: initial codes derived a priori from frameworks and prior studies, applied as a template, with inductive codes added through documented rules when observations do not fit. The three-field code definition (label, definition, how to know it occurs) is a minimal alternative to MacQueen's six fields. Stage 2's pilot on a few documents with a second coder, and stage 6's corroboration step, map onto the thesis' pilot and its check that themes are traceable to coded observations. The paper does not quantify agreement (a supervisor check with no statistic), so it should support the hybrid logic only, not the reliability claims.
\end{papercard}
